\begin{helper}
Fix a history cell, terminal coordinate set \(U\), a red support label \(B\),
a blue label \(J\), and the branch/transcript interface supplied by
\(\noderef{FixedSetHistoryCellBlueBranchTranscriptCells}\).  Assume explicit
released-block data for this fixed pair of labels.  Then there is a residual
assignment mass \(\mu\), released residual blocks
\(\mathcal R_{\beta,\tau}\), and released block masses \(\rho_{\beta,\tau}\)
such that
\[
0\le \mu(A),\qquad \sum_{A\subseteq U}\mu(A)\le1,
\]
each \(\rho_{\beta,\tau}\) is the \(\mu\)-mass of the corresponding released
block, the blocks for distinct realized pairs are disjoint, and the released
cylinder lower bound
\[
p^{|P_\tau|-u_R-u_B}(1-p)^{|A_\tau|-|Z_\tau|}
\le \rho_{\beta,\tau}
\]
holds for every realized pair.
\end{helper}
\begin{proof}
SKETCH:
Use the supplied released-block data exactly as in the red residual partition
node.  Define \(\mu(A)=p^{|A|}(1-p)^{|U|-|A|}\) on assignments
\(A\subseteq U\) and \(0\) outside \(2^U\).  The finite product identity
\(\noderef{FinsetPowersetCylinderMass}\) gives total mass at most \(1\), and
the block hypotheses give the exact formulas for \(\rho_{\beta,\tau}\),
disjointness, and cylinder containment.  The realized-cell geometry from
\(\noderef{FixedSetHistoryCellBlueBranchTranscriptCells}\) rewrites the
contained cylinder mass into the displayed lower bound.
\end{proof}
